Table~\ref{tab:main} gives the main grid (5 seeds, mean $\pm$ std), Fig.~\ref{fig:curves} the same trends, and Table~\ref{tab:paired} the paired differences to CE. At 0\% noise there are no corrupted examples, so the memorisation columns are 0 by definition.

\begin{table*}[t]\centering\caption{Main grid: final-epoch clean test accuracy and memorisation of the corrupted training labels (mean $\pm$ std over 5 seeds; bold = best, underline = second best within a noise rate). Small-loss uses the true noise rate.}\label{tab:main}
\resizebox{0.8\textwidth}{!}{\begin{tabular}{llccccc}
\toprule
Method & Task & test\_acc $\uparrow$ & mem\_rate $\downarrow$ & recover\_rate $\uparrow$ & mem\_gap $\downarrow$ & $n$ \\
\midrule
CE & digits\_noise0 & 0.991 $\pm$ 0.003 & 0.000 $\pm$ 0.000 & 0.000 $\pm$ 0.000 & 0.000 $\pm$ 0.000 & 5 \\
Mixup & digits\_noise0 & \underline{0.993 $\pm$ 0.002} & 0.000 $\pm$ 0.000 & 0.000 $\pm$ 0.000 & 0.000 $\pm$ 0.000 & 5 \\
Label smoothing & digits\_noise0 & \textbf{0.995 $\pm$ 0.001} & 0.000 $\pm$ 0.000 & 0.000 $\pm$ 0.000 & 0.000 $\pm$ 0.000 & 5 \\
\textbf{Small-loss (1 net)} & digits\_noise0 & 0.991 $\pm$ 0.003 & 0.000 $\pm$ 0.000 & 0.000 $\pm$ 0.000 & 0.000 $\pm$ 0.000 & 5 \\
\midrule
CE & digits\_noise20 & 0.897 $\pm$ 0.017 & 1.000 $\pm$ 0.000 & 0.000 $\pm$ 0.000 & 1.000 $\pm$ 0.000 & 5 \\
Mixup & digits\_noise20 & \underline{0.921 $\pm$ 0.013} & \underline{0.919 $\pm$ 0.019} & \underline{0.079 $\pm$ 0.019} & \underline{0.839 $\pm$ 0.038} & 5 \\
Label smoothing & digits\_noise20 & 0.894 $\pm$ 0.012 & 1.000 $\pm$ 0.000 & 0.000 $\pm$ 0.000 & 1.000 $\pm$ 0.000 & 5 \\
\textbf{Small-loss (1 net)} & digits\_noise20 & \textbf{0.943 $\pm$ 0.017} & \textbf{0.419 $\pm$ 0.056} & \textbf{0.550 $\pm$ 0.063} & \textbf{-0.132 $\pm$ 0.119} & 5 \\
\midrule
CE & digits\_noise40 & 0.709 $\pm$ 0.035 & 1.000 $\pm$ 0.000 & 0.000 $\pm$ 0.000 & 1.000 $\pm$ 0.000 & 5 \\
Mixup & digits\_noise40 & \underline{0.762 $\pm$ 0.030} & \underline{0.892 $\pm$ 0.006} & \underline{0.073 $\pm$ 0.016} & \underline{0.819 $\pm$ 0.020} & 5 \\
Label smoothing & digits\_noise40 & 0.668 $\pm$ 0.028 & 0.999 $\pm$ 0.002 & 0.000 $\pm$ 0.000 & 0.999 $\pm$ 0.002 & 5 \\
\textbf{Small-loss (1 net)} & digits\_noise40 & \textbf{0.896 $\pm$ 0.018} & \textbf{0.280 $\pm$ 0.037} & \textbf{0.649 $\pm$ 0.048} & \textbf{-0.370 $\pm$ 0.084} & 5 \\
\midrule
CE & digits\_noise60 & 0.436 $\pm$ 0.039 & 0.995 $\pm$ 0.008 & 0.001 $\pm$ 0.002 & 0.994 $\pm$ 0.010 & 5 \\
Mixup & digits\_noise60 & \underline{0.523 $\pm$ 0.032} & \underline{0.853 $\pm$ 0.049} & \underline{0.078 $\pm$ 0.020} & \underline{0.775 $\pm$ 0.066} & 5 \\
Label smoothing & digits\_noise60 & 0.410 $\pm$ 0.041 & 0.994 $\pm$ 0.007 & 0.001 $\pm$ 0.002 & 0.993 $\pm$ 0.009 & 5 \\
\textbf{Small-loss (1 net)} & digits\_noise60 & \textbf{0.771 $\pm$ 0.035} & \textbf{0.230 $\pm$ 0.025} & \textbf{0.625 $\pm$ 0.035} & \textbf{-0.395 $\pm$ 0.059} & 5 \\
\bottomrule
\end{tabular}
}
\end{table*}

\begin{figure}[t]\centering\includegraphics[width=\columnwidth]{noise_curves.pdf}
\caption{Test accuracy (left) and mem\_rate (right) against the symmetric noise rate; mean $\pm$ std over 5 seeds.}\label{fig:curves}\end{figure}

\begin{table}[t]\centering\caption{Paired differences to CE over the same 5 seeds (system minus CE) and paired $t$-test $p$-value for test accuracy; $n{=}5$, no multiple-comparison correction.}\label{tab:paired}
\resizebox{0.85\columnwidth}{!}{\begin{tabular}{llrrr}
\toprule
Noise & System & $\Delta$acc & $\Delta$mem\_rate & $p_{\rm acc}$ \\
\midrule
20\% & LS & -0.003 & +0.000 & 0.5740 \\
20\% & Mixup & +0.024 & -0.081 & 0.0429 \\
20\% & Small-loss & +0.046 & -0.581 & 0.0008 \\
\midrule
40\% & LS & -0.041 & -0.001 & 0.0141 \\
40\% & Mixup & +0.053 & -0.108 & 0.0003 \\
40\% & Small-loss & +0.187 & -0.720 & $<$0.0001 \\
\midrule
60\% & LS & -0.026 & -0.000 & 0.0948 \\
60\% & Mixup & +0.087 & -0.142 & 0.0018 \\
60\% & Small-loss & +0.335 & -0.765 & $<$0.0001 \\
\bottomrule
\end{tabular}
}
\end{table}

\begin{table}[t]\centering\caption{Mean test accuracy from the logged 5-epoch curves of the main-grid runs (5 seeds): epoch 10, best logged checkpoint (selected on the test split, hence optimistic) and final epoch 60.}\label{tab:curves}
\resizebox{0.9\columnwidth}{!}{\begin{tabular}{llccc}
\toprule
Noise & System & epoch 10 & best ckpt & epoch 60 \\
\midrule
20\% & CE & 0.862 & 0.921 & 0.897 \\
20\% & LS & 0.916 & 0.942 & 0.894 \\
20\% & Mixup & 0.936 & 0.954 & 0.921 \\
20\% & Small-loss & 0.862 & 0.957 & 0.943 \\
\midrule
40\% & CE & 0.763 & 0.805 & 0.709 \\
40\% & LS & 0.739 & 0.854 & 0.668 \\
40\% & Mixup & 0.873 & 0.891 & 0.762 \\
40\% & Small-loss & 0.763 & 0.904 & 0.896 \\
\midrule
60\% & CE & 0.603 & 0.634 & 0.436 \\
60\% & LS & 0.578 & 0.663 & 0.410 \\
60\% & Mixup & 0.655 & 0.728 & 0.523 \\
60\% & Small-loss & 0.603 & 0.779 & 0.771 \\
\bottomrule
\end{tabular}
}
\end{table}

\textbf{Evaluation epoch.} The run logs record test accuracy every 5 epochs. Table~\ref{tab:curves} shows that the final-epoch ranking is specific to end-of-training evaluation: at 40\% noise the best logged checkpoint is 0.805 for CE, 0.854 for LS, 0.891 for mixup and 0.904 for small-loss, so LS exceeds CE and mixup nearly matches small-loss before memorisation sets in, while small-loss is the only system whose final accuracy stays close to its best checkpoint. The best-checkpoint column picks the maximum on the test split and is an upper bound, not a protocol a practitioner could follow.

\textbf{H1 (CE memorises): supported.} CE fits essentially all corrupted labels (mem\_rate 1.000 at 20\% and 40\% noise, 0.995 at 60\%) even though the network has only about 41k parameters and weight decay. Its test accuracy falls from 0.991 at 0\% noise to 0.709 at 40\% and 0.436 at 60\%, i.e.\ by far more than the 0.10 of H1 (at 20\% noise it already falls to 0.897). Final-epoch accuracy is measured with no early stopping, so these numbers reflect the end of training, not the best epoch.

\textbf{H2 (label smoothing gives no material gain): supported, and it is worse than that.} With $\epsilon{=}0.1$, LS is indistinguishable from CE at 20\% (paired $p{=}0.5740$), and is \emph{below} CE at 40\% ($\Delta=-0.041$, $p{=}0.0141$, uncorrected) and at 60\% ($\Delta=-0.026$, $p{=}0.0948$); mem\_rate is unchanged (about 1.0). LS is the best system only at 0\% noise (0.995 vs 0.991 for CE, a difference of 0.4 points over 5 seeds, which we do not interpret). The $\epsilon$ sweep (Table~\ref{tab:sens}) gives test accuracy from 0.662 to 0.693 over $\epsilon\in[0.05,0.6]$ against CE's 0.709, with no value reducing mem\_rate, so within this range the conclusion does not depend on $\epsilon$. This contrasts with Lukasik et al.~\cite{lukasik2020does}, who find label smoothing competitive under label noise; our setting (tiny data, a network that memorises completely, final-epoch evaluation) differs from theirs, and Table~\ref{tab:curves} shows that LS does exceed CE before memorisation sets in (best checkpoint at 40\%: 0.854 vs 0.805), so the disagreement is about end-of-training behaviour.

\textbf{H3 (mixup helps): supported at $\alpha{=}1$, with a caveat on mechanism.} Mixup beats CE at 40\% and 60\% noise clearly and at 20\% only suggestively (the 20\% result rests on an uncorrected $p$ of 0.0429 and one seed goes the other way) ($\Delta=+0.024,+0.053,+0.087$; paired $p$ 0.0429, 0.0003, 0.0018) and lowers mem\_rate (0.919, 0.892, 0.853 against about 1.0). It is nevertheless far from solving the problem: it still memorises 85--92\% of the corrupted labels and recovers the true class for only 7--8\%, so its accuracy gain does not come from avoiding memorisation; we conjecture it comes from the smoothness of the memorised solution, but no experiment here tests that mechanism. Its gain grows with the noise rate, as with the sample-selection rule.

\textbf{H4 (small-loss is best at 40/60\% but loses at 0\%): half supported.} With the true noise rate supplied, small-loss selection is the most accurate system at 20, 40 and 60\% (0.943, 0.896, 0.771; $\Delta=+0.187$ and $+0.335$ over CE at 40 and 60\%, $p<0.0001$) and the only one with a negative mem\_gap (it classifies more corrupted training images by their true class than by the wrong one: $-0.370$ at 40\% and $-0.395$ at 60\%). The second clause is not testable as stated: at 0\% noise the forget rate is 0, so the system is numerically identical to CE (same seeds, identical accuracy 0.991), and it does not lose to CE. The cost of selection should instead be read from where it fails: at 60\% noise it still ends at 0.771, well below its own 0\% accuracy of 0.991, and at 20\% noise it fits 42\% of the corrupted labels (mem\_rate 0.419) because the forget rate ramps only after a warm-up in which the network is trained on all examples.

\textbf{Failure cases.} (i) At 60\% noise every system is far from the clean-label accuracy: CE 0.436, LS 0.410, mixup 0.523, small-loss 0.771. (ii) Small-loss selection relies on knowing the noise rate; this is not available in practice and Sec.~\ref{sec:ablations} shows what happens when it is wrong. (iii) The per-seed spread of mem\_gap for small-loss at 20\% noise is large ($-0.132\pm0.119$): the same recipe sometimes ends with the memorised and sometimes with the clean label dominant.
